\section{Introduction}
\label{sec:introduction}

% ------------------------------------------------------------
% Paragraph 1 -- Background: microservices + distributed tracing
% ------------------------------------------------------------
Microservice architectures decompose large applications into collections of independently deployed services that communicate through networked requests  \cite{dragoni2017microservices}. Although this design improves scalability, modularity, and deployment flexibility, it also complicates performance diagnosis because a single user request may traverse many services and generate a distributed execution path. Distributed tracing systems record this execution as a trace composed of spans, where each span captures a service operation, timing information, and parent-child relationships \cite{sigelman2010dapper}. As microservice systems scale, manual inspection of traces becomes impractical, creating a need for automated methods that can identify traces whose execution deviates from expected behavior and guide operators toward likely sources of latency or failure.

% ------------------------------------------------------------
% Paragraph 2 -- Why trace anomaly detection matters
% ------------------------------------------------------------
At the scale of modern microservice deployment, distributed tracing
frameworks continuously collect large volumes of execution traces,
making manual inspection challenging for human operators. Early
industrial approaches often relied on manually maintained rules,
including service-specific return-code rules, to identify abnormal
executions. As services evolve, however, these rules require continual
maintenance because the meaning and usage of service-specific return
codes may change over time, reducing their long-term effectiveness
\cite{liu2020traceanomaly}. These challenges have motivated
learning-based trace anomaly detection methods that model the
characteristics of normal trace executions and automatically identify
traces whose behavior deviates from the learned patterns, allowing
operators to prioritize investigation without continually redesigning
rule sets. 
% ------------------------------------------------------------
% Paragraph 3 -- Existing methods and the STV representation
% ------------------------------------------------------------
Learning-based trace anomaly detection depends heavily on how
distributed traces are represented for machine learning. Previous
studies have explored execution sequences, graph-based models,
learned embeddings, and fixed-dimensional vector representations \cite{liu2020traceanomaly,zhang2022deeptralog,xie2023tracevae}. Among these, the Service Trace Vector (STV), introduced in
TraceAnomaly \cite{liu2020traceanomaly}, transforms variable-length
distributed traces into fixed-dimensional feature vectors suitable for
machine learning. In an STV, each feature corresponds to a specific
service reached through a particular call path, while its value
represents a latency-related characteristic associated with that
service path (Section~\ref{sec:background}). Because every feature has
a direct correspondence to an actual service path, STVs preserve
service-level interpretability while providing a compact numerical
representation of distributed traces. Furthermore, because STVs are
fixed-dimensional feature vectors, they can be analyzed using a wide
range of conventional statistical and machine learning techniques,
making them a flexible representation for trace anomaly detection.
However, the effectiveness of such approaches depends not only on the
trace representation itself, but also on how anomaly scores are
derived from the resulting feature vectors.Existing STV- and graph-based approaches have demonstrated strong anomaly detection performance across several benchmark microservice datasets. However, comparatively little attention has been given to understanding whether the anomaly scores produced by these methods are influenced by intrinsic structural properties of traces, such as execution size or call-path complexity.


% ------------------------------------------------------------
% Paragraph 4 -- The research gap: complexity bias
% ------------------------------------------------------------
We identify a systematic weakness in STV-based anomaly scoring: when an STV
is scored using a generic unsupervised detector trained directly on raw,
unnormalized feature values, the resulting anomaly score can become
strongly influenced by trace \emph{size}, independent of trace
\emph{health}. To better understand this phenomenon, we analyze real, non-injected
test traces from three configuration variants of the
\texttt{ts-order-service} microservice in the TrainTicket benchmark \cite{zhou2021trainticket}.
Across these datasets, the Pearson correlation between baseline
STV~+~Isolation Forest anomaly scores and the number of distinct
services ranges from $r = 0.968$ to $0.982$, while the correlation with
the number of spans ranges from $r = 0.854$ to $0.981$
(Table~\ref{tab:complexity-bias-baseline}). Because these measurements are obtained from traces without injected
anomalies, the observed correlations cannot be explained by successful
detection of synthetic faults. Instead, they indicate that the structural characteristics of a trace can
substantially influence anomaly scores even when the observed executions
are operationally normal. This suggests that anomaly scoring based solely on raw STV feature
values may find it more difficult to distinguish genuinely abnormal
executions from structurally complex but otherwise normal traces.
% ------------------------------------------------------------
% Paragraph 5 -- Proposed method
% ------------------------------------------------------------
We propose \textbf{WR-STV} (Weighted Residual Service Trace Vector), which
replaces the raw STV value at each feature with a standardized residual
against that feature's learned training-time mean and variance, clips the
residual for numerical stability, weights it by a log-frequency reliability
term that down-weights rarely observed paths, and scores the trace by the
mean of its top-$k$ weighted absolute residuals. This keeps the
representation's per-service-path interpretability while reducing its dependence on trace size (Section~\ref{sec:methodology}).

% ------------------------------------------------------------
% Contributions
% ------------------------------------------------------------
The main contributions of this paper are:

\begin{itemize}
    \item \textbf{An empirical demonstration of trace-complexity bias} in
    STV~+~Isolation Forest anomaly scoring, measured directly on real,
    non-injected test traces from three TrainTicket configurations
    ($r \approx 0.85$--$0.98$ with baseline score; reduced to
    $r \approx 0.24$--$0.30$ under WR-STV).

    \item \textbf{Weighted Residual Service Trace Vectors}: a lightweight residual-standardization and
    reliability-weighting scheme for STV scoring, together with an ablation
 study isolating how much of its improvement comes from standardization versus
    weighting.

    \item \textbf{A controlled, circularity-checked evaluation}: evaluate WR-STV using latency
    anomalies injected into both the largest-duration visible span and a randomly chosen visible span (robustness check), injection-strength sensitivity, classical  anomaly detection baselines, and cross-dataset transfer across all six directed train/test pairs. Showing that WR-STV's detection improvement is not an artifact of always targeting the feature to which its scoring function is most sensitive.

   
    \item \textbf{Quantitative interpretability and failure analysis}:
    top-$k$ explanation overlap against known injected paths, and a
    taxonomy of the conditions under which WR-STV fails to rank an injected
    trace above its original.
\end{itemize}

